\documentclass[12pt]{article}
\usepackage[utf8]{inputenc}
\usepackage[T1]{fontenc}
\usepackage[margin=1in]{geometry}
\usepackage{setspace}
\usepackage{xcolor}
\usepackage[colorlinks=true, urlcolor=blue!60!black, linkcolor=blue!60!black]{hyperref}
\setlength{\parindent}{0em}
\setlength{\parskip}{0.6em}
\renewcommand{\baselinestretch}{1.0}
\begin{document}
\begin{center}
{\LARGE \textit{New York State Rifle \& Pistol Association, Inc. v. Bruen}, 597 U.S. 1 (2022)\par}
\vspace{0.5em}
{\large Case Brief --- Two-Page Edition\par}
\end{center}

\section{Memory Jogger}
Law-abiding citizens have a right to carry handguns in public for self-defense, and gun laws are judged by text and history, not means-end scrutiny (\href{https://www.law.cornell.edu/supremecourt/text/20-843}{597 U.S. 1}).

\section{Detailed Case Facts}
New York required applicants for an unrestricted concealed-carry license to show ``proper cause,'' meaning a special need for self-protection beyond that of the general community; open carry was banned. Robert Nash and Brandon Koch, law-abiding members of the New York State Rifle \& Pistol Association, were limited to hunting and target-shooting licenses (Koch could also carry to and from work).

\section{Procedural History}
Bound by \textit{Kachalsky}, the district court dismissed and the Second Circuit affirmed; the Supreme Court reversed.

\section{Judicial Votes}
6--3. Thomas wrote for Roberts, Alito, Gorsuch, Kavanaugh, and Barrett; Alito, Kavanaugh (with Roberts), and Barrett concurred; Breyer dissented, joined by Sotomayor and Kagan.

\section{Holding}
When the plain text covers conduct, the government must show its regulation is consistent with the Nation's historical tradition, reasoning by analogy for new problems. The proper-cause requirement fails.

\section{Analysis of Opinions}
Thomas called the two-step test ``one step too many'' and found New York's historical evidence wanting: English and colonial laws barred carrying to terrify, surety laws presumed a right to carry, and Western territorial bans were outliers. Kavanaugh stressed that objective shall-issue licensing and longstanding prohibitions survive; Barrett flagged the 1791-versus-1868 question. Breyer, citing gun-violence data, said the history-only test lets judges ``pick their friends out of history's crowd.''

\section{Examples: Future Applications}
\textbf{Same-Side:} a discretionary ``good moral character'' test without objective criteria is invalid. \textbf{Opposite-Side:} objective shall-issue licensing with background checks and training is valid. \textbf{Fence-Sitter:} a default ban on carry on private property open to the public absent the owner's consent.

\section{Critique}
Originalists such as Baude and Leider fault the hunt for analogues; Blocher, Ruben, and Charles say ``originalism-by-analogy'' confuses lower courts, a problem \textit{Rahimi} tried to address.

\section{Key Quotations}
\begin{quote}
``When the Second Amendment's plain text covers an individual's conduct, the Constitution presumptively protects that conduct. The government must then justify its regulation by demonstrating that it is consistent with the Nation's historical tradition of firearm regulation. Only then may a court conclude that the individual's conduct falls outside the Second Amendment's `unqualified command.' ''
\end{quote}
\noindent --- Justice Thomas, Opinion of the Court, 597 U.S. 1.

\end{document}
